\begin{table}[htbp]
\centering
\caption{Prospectively specified numerical-validation design}
\label{tab:validation-design}
\begin{tabularx}{\textwidth}{p{0.19\textwidth} X X}
\toprule
Domain & Prospectively specified design & Purpose \\
\midrule
Scientific units & 25 latent equivalence classes: 17 primary and 8 diagnostic. Fifty deterministic common-monotone-transform rows were used only for invariance audits and did not increase the formal denominator. & Primary classes determine confirmatory acceptance; diagnostic classes and transformed rows probe boundary behavior and implementation invariance. \\
Simulation structure & Three families—reference-only, evaluation-only, and combined—with 25 jobs per family (75 jobs total). Reference-bank and evaluation-sample sizes were 250, 500, 1,000, 3,000, and 10,000; TESS was evaluated at alpha=0.01 and 0.05. & Separates reference-threshold error, finite-n evaluation error, and their joint propagation. \\
Random-number contract & Master seed 20260804, independent reference and evaluation streams, nested prefixes within each family, atomic batch writes, and deterministic resumption. & Ensures exact replay and interruption-safe execution. \\
Replication and stopping & Primary classes: 2,000–20,000 replicates; diagnostic classes: 1,000–5,000 replicates; batch size 250. Stopping required MCSE/(1+|tau|) <=0.03, where tau was the locked population target; no observed Monte Carlo mean entered the stopping rule. & Controls Monte Carlo resolution without effect-dependent stopping. \\
Reference approximation & At B=10,000, median normalized residual <=0.15 and 90th percentile <=0.40; median improvement from B=500 to B=10,000 >=40\%; generalized correction improved at least 75\% of primary classes relative to the prespecified baseline. & Assesses finite-sample usefulness of the generalized reference correction. \\
Combined policy approximation & At (B,n)=(10,000,10,000), median normalized residual <=0.15 and 90th percentile <=0.40. & Assesses the joint reference/evaluation bias prediction, including the B-by-n interaction term. \\
Combined TESS approximation & At (B,n)=(10,000,10,000), separately for each alpha, median normalized residual <=0.25 and 90th percentile <=0.60; any nonfinite primary record constituted formal failure. & Assesses nonlinear propagation to TESS with explicit boundary-status accounting. \\
Exact evaluation identity & Eighty-five two-sided Bonferroni comparisons at familywise alpha=0.01 (z*=3.85098), targeting E(delta\_hat)=(1-1/n)Delta\_pi. & Provides a prespecified fatal implementation check. \\
Formal decision rule & All scientific criteria and fatal checks had to pass. No more than 10\% of primary classes could be precision limited; with 17 primary classes, the operational maximum was one class. Post-hoc class dropping was prohibited. & Preserves the predeclared confirmatory denominator and prevents a pass driven by inadequate Monte Carlo precision. \\
\bottomrule
\end{tabularx}
\begin{minipage}{\textwidth}
\footnotesize
\textsuperscript{1} B denotes the reference-bank size, n the evaluation-sample size, alpha the TESS level, and tau the locked population target used only for the Monte Carlo stopping scale.\\
\textsuperscript{2} Normalized residuals use the prospectively specified simultaneous Monte Carlo uncertainty adjustment; raw-residual diagnostics are reported with the numerical results.\\
\textsuperscript{3} Detailed seed, transformation-audit, and role-specific replication specifications are provided in Supplementary Table S1.\\
\end{minipage}
\end{table}
